\section{Discussion}
\label{sec:discussion}

This section interprets the results with respect to the three research
questions (Section~\ref{sec:rq}), examines the label-format issue in
broader context (Section~\ref{sec:disc-label}), discusses the deployment
pipeline (Section~\ref{sec:disc-deploy}), and states the limitations of
the current work (Section~\ref{sec:limitations}).

% -------------------------------------------------------------------------
\subsection{Research Questions Revisited}
\label{sec:rq}
% -------------------------------------------------------------------------

\textbf{RQ1: How accurate and fast is YOLO26m on this dataset?}

The deployed YOLO26m checkpoint achieves mAP$_{50}$~=~0.784 and
mAP$_{50:95}$~=~0.594 on the test split (Table~\ref{tab:overall}).
These figures place the model in a competitive range for applied vehicle
detection, though direct comparison with published studies is confounded
by dataset differences.
On the COCO benchmark, YOLO26m is reported to achieve mAP$_{50:95}$ in
the mid-to-high 50s \cite{yolo262023}; the value of 0.594 obtained here
suggests that the three-class vehicle dataset is of similar or slightly
lower difficulty than the full 80-class COCO evaluation for the classes
in question.

Precision (0.845) substantially exceeds recall (0.752), implying that
the model is conservative — when it makes a prediction, it is usually
correct, but it misses a non-trivial fraction of vehicles.
The TIDE analysis confirms this: missed detections (78) outnumber
false positives from background (22) by more than three to one.
For traffic-counting applications, false negatives are typically more
harmful than false positives because under-counting leads to incorrect
density estimates.
This imbalance could be addressed by lowering the confidence threshold
below 0.40 at the cost of increased background detections.

A fair comparison with YOLOv8m, YOLO11m, and RT-DETR-l is [DATA
DIBUTUHKAN] because the \texttt{train} stage was not executed.
The architectural expectation is that YOLO26's NMS-free inference will
show a small latency advantage over NMS-using baselines, particularly
under high-density scenes where NMS operates on many candidate boxes.

\textbf{RQ2: How common are mixed label files, and what is their effect?}

Three of 9\,218 training label files contained mixed bounding-box and
polygon rows (0.033\,\% of files).
All three belonged to the \textit{bus} class, which is consistent with
the hypothesis that bus outlines were annotated as polygons by a subset of
contributors while other classes used the standard bounding-box format.

Because the validation and test label files had no mixed rows, the
evaluation metrics in Table~\ref{tab:overall} are identical for original
and cleaned labels.
The impact of the mixed-format bug therefore manifests at \emph{training
time}: three label files in the training split were silently misread by
Ultralytics~8.4.53, corrupting the box coordinates for every other
bounding-box annotation in those files.
Measuring the accuracy change requires re-training with and without the
fix and is [DATA DIBUTUHKAN] from the \texttt{train} stage.

The practical implication is significant even before the numbers are
available.
The bug produces no error message and no visible warning in the training
log, making it invisible to practitioners who do not audit their label
files.
Our audit script (\texttt{stage\_audit} in \texttt{trace\_experiments.py})
takes seconds to run on a dataset of this size and reliably identifies all
mixed files.
We recommend running such an audit before any training run on
community-sourced YOLO-format datasets.

\textbf{RQ3: Is the deployment pipeline suitable for consumer hardware?}

At 99.1\,ms per image on CPU (Apple M4 Pro, 14-core ARM), the model
processes approximately 10 images per second, which is below the threshold
for smooth real-time video but adequate for frame-sampled analysis
(e.g., one frame per second from a fixed traffic camera) or offline batch
processing.
The Streamlit application introduces additional latency for image loading,
Streamlit rendering, and optional n8n webhook calls
(measured separately; MPS latency is [DATA DIBUTUHKAN]).

Pre-processing (0.62\,ms) and post-processing (0.08\,ms) are negligible
compared with the neural-network forward pass (96.31\,ms), indicating
that YOLO26's NMS-free design has fully eliminated the NMS overhead in the
post-processing stage.
For comparison, a standard YOLO detector with NMS typically adds 1–5\,ms
of post-processing latency, depending on the number of candidate boxes
\cite{TODO-nms-free}.
On Apple MPS the forward pass is expected to run significantly faster;
informal testing suggests sub-50\,ms per image, which would bring the
throughput above 20\,FPS.

% -------------------------------------------------------------------------
\subsection{Label-Format Issues in Community Datasets}
\label{sec:disc-label}
% -------------------------------------------------------------------------

The mixed bounding-box/polygon issue encountered here likely occurs in
other datasets on Roboflow Universe, particularly those assembled by
merging projects with different annotation policies.
Roboflow's polygon-to-YOLO export converts polygons to polygons rows
(many tokens) and bounding boxes to bbox rows (five tokens) correctly
within a single project.
The problem arises when datasets from multiple projects are merged and
some contributors used polygon tools while others used rectangle tools.

The scope of this problem across the full Roboflow Universe catalogue is
unknown.
The \texttt{trace\_novelty.py} harness in this repository includes a
planned stage that downloads a sample of public vehicle datasets and
runs the same audit to estimate the prevalence of mixed files.
Until that audit is complete, practitioners should treat community
YOLO-format datasets as potentially containing mixed files and run
a lightweight token-count check before training.

% -------------------------------------------------------------------------
\subsection{Deployment Pipeline}
\label{sec:disc-deploy}
% -------------------------------------------------------------------------

The VehicleScope~AI Streamlit application demonstrates that a
state-of-the-art NMS-free detector can be deployed for practical use
without specialised infrastructure.
The n8n webhook and Google Sheets logging layer allows non-technical
users to access counts via a familiar spreadsheet interface without
writing any code.
The Telegram alert mechanism provides a low-latency notification channel
for high-occupancy events (more than 10 vehicles per frame).

One limitation of the current deployment is that it processes individual
uploaded images rather than a live video stream.
Extending the pipeline to ingest a real-time RTSP stream would require a
frame-capture thread and either a queue-based architecture or frame
dropping to maintain acceptable latency.
At 10.1\,FPS on CPU the system could process one out of every three
frames of a 30\,FPS feed without buffering; MPS acceleration would
improve this ratio.

% -------------------------------------------------------------------------
\subsection{Limitations}
\label{sec:limitations}
% -------------------------------------------------------------------------

The following limitations should be considered when interpreting the
results.

\begin{enumerate}

  \item \textbf{Single checkpoint, no re-training comparison.}
    All quantitative results are from the single pre-trained \texttt{best.pt}
    checkpoint.
    The baseline comparison (YOLO26m vs.\ YOLOv8m, YOLO11m, RT-DETR-l)
    and the augmentation ablation have not yet been executed.
    Claims about YOLO26's superiority over baselines are therefore
    conditional on future experiments.

  \item \textbf{No cross-domain evaluation.}
    The dataset originates from a single Roboflow project whose geographic
    provenance is not fully documented.
    How well the trained model generalises to footage from a different
    city or camera network is unknown.
    The planned Surabaya traffic dataset (contribution~C3,
    \texttt{trace\_novelty.py}) would address this gap.

  \item \textbf{CPU-only latency measurement.}
    The MPS latency run did not complete successfully in the current
    experiment set.
    The CPU figure (99.1\,ms) over-represents the deployment latency on
    the target Apple M4 Pro hardware.

  \item \textbf{Label provenance.}
    The dataset does not publish per-image metadata such as recording
    location, date, or camera model.
    It is therefore not possible to control for confounds such as seasonal
    lighting variation or camera-specific artefacts.

  \item \textbf{Class definition ambiguity.}
    The boundary between \textit{car} and \textit{van} is defined by
    annotator judgement and is likely inconsistent across contributors.
    This ambiguity is a structural property of the dataset rather than a
    failure of the detection model, but it limits the ceiling of achievable
    accuracy for the \textit{van} class.

  \item \textbf{Single random seed.}
    The deployed checkpoint was trained with one seed.
    Metric variance across seeds is unknown; the planned three-seed
    protocol (seeds 0, 1, 2) will quantify this.

\end{enumerate}
